\section{Related Work}\label{sec:RelatedWork}
    Remote sensing change detection is closely related to visual representation learning, bi-temporal feature comparison, pseudo-change suppression, and dense prediction. This section reviews representative studies from these four aspects.

    \subsection{Visual Representation Learning for Remote Sensing}\label{sec:RelatedWork:Vision}
        Remote sensing imagery contains complex spatial structures, spectral responses, and scale variations, which makes robust visual representation learning essential for downstream interpretation. Ji \etal proposed PASSNet\cite{passnet} with a patch attention module for hyperspectral image classification. Han \etal studied optical remote sensing object detection using weakly supervised learning and high-level feature learning\cite{weakly-supervised}. Wang \etal designed FSoD-Net\cite{Full-scale} for full-scale detection in optical remote sensing imagery. Liu \etal investigated unsupervised contextual anomaly detection for remotely sensed data\cite{unsupervised}. Xiao \etal introduced a diffusion-based super-resolution model for remote sensing images\cite{diffusion}. These studies improve remote sensing representation from classification, detection, anomaly modeling, and restoration perspectives, but they mainly operate on single images. In contrast, change detection requires discriminative representation as well as explicit comparison between two temporal observations.

    \subsection{Bi-Temporal Remote Sensing Change Detection}\label{sec:RelatedWork:ChangeDetection}
        Deep change detection methods usually learn bi-temporal representations with siamese or dual-stream architectures. Daudt \etal introduced fully convolutional siamese networks as an early deep comparison framework\cite{fully-rs}. Fang \etal further improved siamese modeling through dense connections in SNUNet-CD\cite{Siamese-network}. Long \etal proposed SASiamNet\cite{Self-adaptive-Siamese-network} with self-adaptive feature interaction in a siamese framework. Peng \etal combined attention mechanisms with image differencing for change detection\cite{attention-sample}. Wang \etal introduced attention-based deep supervision in ADS-Net\cite{deeply-supervised-network}. Gu \etal focused on full-scale difference feature fusion in FDFF-Net\cite{full-scale-difference-feature-fusion-network}. Wang \etal proposed a high-resolution feature difference attention network for building change detection\cite{feature-difference-attention-network}. Chen \etal introduced transformer-based modeling for remote sensing change detection\cite{transformers-rs}. Ma \etal presented EATDer\cite{ma2023eatder} with edge-assisted adaptive transformer detection. Li \etal further combined semantic flow and edge-aware refinement in SFEARNet\cite{sfearnet} to improve efficient change reasoning. Although these methods have advanced bi-temporal fusion and boundary localization, many of them still treat visual discrepancy as the main evidence of change, which may lead to false responses in pseudo-change regions.
        
    \IEEEpubidadjcol
    
    \subsection{Pseudo-Change Suppression in Change Detection}\label{sec:RelatedWork:PseudoChange}
        Pseudo-change refers to visually apparent differences that are not annotated as semantic changes. It is commonly caused by illumination variation, seasonal transition, sensor inconsistency, atmospheric interference, or local imaging conditions. Zhang and Xue proposed SIAM-CDNET\cite{optimized-edge-detection} to optimize edge detection and mitigate pseudo changes. Cheng \etal introduced a content-cleansing framework to remove unreliable cues before change inference\cite{harmony}. Chen \etal formulated building change detection as an XOR-style binary relation problem\cite{buildings}. Wang \etal proposed FIBTNet\cite{Feature-Interactive} to enhance feature interaction between bi-temporal observations. These methods share the view that pseudo-change cannot be fully handled by simple appearance differencing and requires stronger semantic or structural constraints.

        Semi-supervised and contrastive learning methods also provide useful insights for pseudo-change suppression. Sun \etal explored semi-supervised change detection with siamese networks and graph attention\cite{semisanet}. Kondmann \etal proposed SemiSiROC\cite{SemiSiROC} with an unsupervised teacher model. Wang \etal proposed reliable contrastive learning to separate changed and unchanged representations\cite{contrastive-learning}. Wang \etal introduced STCRNet\cite{self-training} with self-training and consistency regularization. These methods improve robustness by exploiting unlabeled data, consistency constraints, or representation separation. However, the core challenge of pseudo-change remains the inconsistency between visual difference and annotation semantics. Therefore, an annotation-consistent supervision strategy is still needed to guide the model toward benchmark-defined changes rather than all observable differences.

    \subsection{Dense Prediction and Semantic Segmentation in Remote Sensing}\label{sec:RelatedWork:DensePrediction}
        Dense prediction methods provide important design references for pixel-wise change detection, especially in foreground modeling and boundary delineation. As a semantic-segmentation reference, Xu \etal proposed RSSFormer\cite{rssformer} to enhance foreground saliency for remote sensing land-cover segmentation. Tang \etal developed WNet\cite{wnet} with a hierarchical W-shaped structure for remote-sensing image change detection. Yuan \etal improved multi-scale feature modeling through the global filter in the frequency domain\cite{multi-scale}. Lu \etal proposed a multi-scale feature progressive fusion network for change detection\cite{multi}. Ma \etal introduced edge-assisted adaptive transformer detection\cite{ma2023eatder}. These methods show that accurate dense prediction benefits from foreground saliency, multi-scale context, and boundary-aware refinement. Nevertheless, semantic segmentation is usually single-temporal, whereas change detection must compare two temporal observations and distinguish semantic change from pseudo-change. Thus, segmentation-related techniques are useful components, but they do not by themselves solve annotation-consistent bi-temporal change detection.
